\documentclass[11pt,a4paper]{article}

\usepackage[margin=2.5cm]{geometry}
\usepackage{fontspec}
\usepackage{amsmath,amssymb,amsthm,mathtools}
\usepackage{booktabs}
\usepackage{microtype}
\usepackage[hidelinks]{hyperref}

\setmainfont{Latin Modern Roman}
\setmonofont{Latin Modern Mono}[Scale=MatchLowercase]

\newtheorem{theorem}{Theorem}[section]
\newtheorem{proposition}[theorem]{Proposition}
\newtheorem{lemma}[theorem]{Lemma}

\newcommand{\LG}{\mathrm{LG2N}}
\newcommand{\fl}[1]{\mathrm{fl}_{#1}}
\newcommand{\emax}{e_{\max}}
\newcommand{\mlim}{m_{\mathrm{lim}}}
\newcommand{\ofs}{\mathrm{ofs}}

\numberwithin{equation}{section}

\title{\textbf{The table-free latitudinal march}}
\author{}
\date{}

\begin{document}
\maketitle
\vspace{-2.2em}

\begin{center}
{\small Lean~4 / Mathlib proof of record: \texttt{formal/lean/} (v4.34.0-rc2).}
\end{center}

\medskip
\noindent\textbf{Abstract.}
We record the identities used by the table-free latitudinal march.  For spin
$s\ge 0$, order $m\ge 0$ and $x=\cos\theta$, the Wigner function is the Jacobi
polynomial
\begin{equation*}
d^\ell_{m,-s}(\theta)
  =(-1)^{m+s}\,
   2^{\varepsilon_m\,\LG(\ell,m)}\,
   \bigl(\sin\tfrac{\theta}{2}\bigr)^{m+s}
   \bigl(\cos\tfrac{\theta}{2}\bigr)^{\lvert m-s\rvert}
   P_{\ell-M}^{(m+s,\lvert m-s\rvert)}(x),
\end{equation*}
with $M=\max(m,s)$ and $\varepsilon_m=\mathrm{sign}(m-s)$ (equal to $+1$ at
$m=s$).  The kernel marches this polynomial by the three-term recurrence,
carries an exact binary exponent, emits a binade-split partial, and may skip
polar tiles by ducc0's $\mlim$ cutoff.  Every identity below is proved in
Lean; only the WKB estimate of \S\ref{sec:polar} is not.

\tableofcontents

\section{Introduction}

The transform implemented by \texttt{gmaster/\_spin\_march\_pallas.py} is, for
each order $m$ and each of four real channels $r_c(\theta)$,
\begin{equation}
\label{eq:A}
A_c(\ell,m)
  =\sum_{\theta} d^\ell_{m,-s}(\theta)\, r_c(\theta),
  \qquad
  \ell=\max(m,s),\dots,L-1,
\end{equation}
where $L=\ell_{\max}+1$.  The kernel never stores associated Legendre tables.
It evaluates $d^\ell_{m,-s}$ from a Jacobi recurrence in compensated float32,
with an integer exponent lane, a float32 emit, and an optional polar skip.

This note is the human-readable statement.  The proof of record is the Lean~4
project \texttt{GMasterMarch}; Appendix~\ref{app:lean} names the theorems.
The project has no \texttt{sorry}: DLMF~18.9.2 for the explicit Jacobi sum
is proved in \texttt{JacobiSum.jacobi\_three\_term} (a telescoping certificate
over the summands, \S\ref{sec:rec}), the identification of the Jacobi closed
form with Wigner's explicit sum in \texttt{WignerJacobi.wignerD\_eq\_jacobiForm}
(\S\ref{sec:closed}), and the spin-$2$ mirror identity in
\texttt{Mirror.mirror\_identity} (\S\ref{sec:fold}).  The WKB estimate in
\S\ref{sec:polar} is asymptotic and has no Lean counterpart.

The argument is organised as follows.  \S\ref{sec:closed} substitutes
$m'\mapsto m$, $m\mapsto -s$ into Edmonds's Jacobi form of the small-$d$
matrix.  \S\ref{sec:rec} rearranges DLMF~18.9.2 into the marched coefficients
and proves that $\alpha^2-\beta^2=4ms$.  \S\ref{sec:carry}--\S\ref{sec:emit}
treat the exponent lane and the binade-split emit.  \S\ref{sec:fold} records
the hemisphere fold.  \S\ref{sec:wall} converts the live-degree count into a
hardware lower bound.  \S\ref{sec:polar} identifies ducc0's $\mlim$ with the
larger root of its quadratic.

\section{Closed form}
\label{sec:closed}

Write $\alpha=m+s$, $\beta=\lvert m-s\rvert$, $M=\max(m,s)$ and $n=\ell-M$.

\begin{theorem}
\label{thm:closed}
\begin{align}
d^\ell_{m,-s}(\theta)
  &=(-1)^{m+s}\,
    2^{\varepsilon_m\,\LG(\ell,m)}\,
    \Bigl(\sin\tfrac{\theta}{2}\Bigr)^\alpha
    \Bigl(\cos\tfrac{\theta}{2}\Bigr)^\beta
    P_n^{(\alpha,\beta)}(x),
\label{eq:d}
\\
\LG(\ell,m)
  &=\log_2\sqrt{\frac{(\ell+m)!\,(\ell-m)!}{(\ell-s)!\,(\ell+s)!}},
  \qquad
  \varepsilon_m
  =
  \begin{cases}
    +1,& m\ge s,\\
    -1,& m<s.
  \end{cases}
\label{eq:LG}
\end{align}
\end{theorem}

\begin{proof}
Edmonds (1957, eq.~4.1.23) writes, for $d^j_{m'm}$ with
$a=\lvert m-m'\rvert$, $b=\lvert m+m'\rvert$ and $n=j-(a+b)/2$,
\begin{equation}
\label{eq:edmonds}
d^j_{m'm}(\beta)
  =\xi\,
   \sqrt{\frac{n!\,(n+a+b)!}{(n+a)!\,(n+b)!}}\,
   \Bigl(\sin\tfrac{\beta}{2}\Bigr)^a
   \Bigl(\cos\tfrac{\beta}{2}\Bigr)^b
   P_n^{(a,b)}(\cos\beta),
\end{equation}
where $\xi=1$ if $m\ge m'$ and $\xi=(-1)^{m-m'}$ otherwise.

Substitute $m'\mapsto m$ and $m\mapsto -s$.  Then $a=\alpha$, $b=\beta$,
and $a+b=2M$, so $n=\ell-M$.  Unless $m=s=0$ one has $-s<m$, hence
$\xi=(-1)^{-s-m}=(-1)^{m+s}$.  The factorial ratio splits on the two
branches of $M$:
\begin{itemize}
\item if $m\ge s$, then $n=\ell-m$, $n+a+b=\ell+m$, $n+a=\ell+s$,
$n+b=\ell-s$, so the ratio is $2^{+\LG}$;
\item if $m<s$, then $n=\ell-s$, $n+a+b=\ell+s$, $n+a=\ell+m$,
$n+b=\ell-m$, so the ratio is $2^{-\LG}$.
\end{itemize}
This is the $\varepsilon_m$ flip in~\eqref{eq:LG}.
\end{proof}

The kernel writes $(-1)^m$ for $(-1)^{m+s}$.  The two agree for even spin
($0$ and $2$, the only spins the march serves).  A spin-$1$ or spin-$3$
march must restore $(-1)^s$.  The closed form \emph{is} Wigner's function:
with $d^j_{m,-s}$ written as Wigner's explicit sum (Wigner 1931; Sakurai
3.8.33), both sides are sums over $0\le k\le j-M$, the powers of
$\sin\frac{\theta}{2}$ and $\cos\frac{\theta}{2}$ match term by term
($\tfrac{x-1}{2}=-\sin^2\frac{\theta}{2}$, $\tfrac{x+1}{2}=\cos^2\frac{\theta}{2}$),
and the coefficients agree by
$\binom{j+s}{j-m-k}\binom{j-s}{k}(j-s-k)!\,(j-m-k)!\,(k+m+s)!\,k!=(j+s)!\,(j-s)!$
for $m\ge s$ (and $(j+m)!\,(j-m)!$ for $m<s$), which is where the
$\varepsilon_m$ flip comes from.  Numerically, spins $0$--$3$, $\ell\le 12$,
four colatitudes: worst relative difference $6\times 10^{-13}$.

Lean: \texttt{Indices.alpha\_add\_beta}, \texttt{alpha\_sub\_beta},
\texttt{alpha\_sq\_sub\_beta\_sq}, \texttt{jratio\_of\_ge},
\texttt{jratio\_of\_lt}, \texttt{sign\_of\_even\_spin},
\texttt{sign\_of\_odd\_spin}; \texttt{WignerJacobi.wignerD} (Wigner's sum),
\texttt{jacobiForm}, \texttt{coeff\_ge}, \texttt{coeff\_lt},
\texttt{wignerD\_eq\_jacobiForm}.

\section{Recurrence}
\label{sec:rec}

DLMF~18.9.2 for $P_n^{(\alpha,\beta)}$, with $S=2n+\alpha+\beta$, is
\begin{equation}
\label{eq:dlmf}
2n\,(n+\alpha+\beta)\,(S-2)\,P_n
  =(S-1)\bigl[S(S-2)x+(\alpha^2-\beta^2)\bigr]P_{n-1}
   -2(n+\alpha-1)(n+\beta-1)S\,P_{n-2}.
\end{equation}
Write $\mathrm{den}=2n(n+\alpha+\beta)(S-2)$.  Then
$v_\ell=(c_1 x+c_0)v_{\ell-1}-c_b v_{\ell-2}$ with
\begin{equation}
\label{eq:c}
c_1=\frac{(S-1)S(S-2)}{\mathrm{den}},
\quad
c_0=\frac{(S-1)(\alpha^2-\beta^2)}{\mathrm{den}},
\quad
c_b=\frac{2(n+\alpha-1)(n+\beta-1)S}{\mathrm{den}}.
\end{equation}

\begin{lemma}
\label{lem:c0}
$\alpha^2-\beta^2=4ms$ for every $m,s\ge 0$.
\end{lemma}

\begin{proof}
$\alpha-\beta=2\min(m,s)$ and $\alpha+\beta=2\max(m,s)$.
\end{proof}

Thus no coefficient in~\eqref{eq:c} is a difference of large numbers.
A float32 (high, low) split therefore carries them to $2^{-48}$.

Seeds: $P_0=1$, and DLMF~18.5.7 at $n=1$ gives
\begin{equation}
\label{eq:P1}
P_1^{(\alpha,\beta)}(x)
  =\tfrac{\alpha+\beta+2}{2}\,x+\tfrac{\alpha-\beta}{2}.
\end{equation}
For $s=2$ and $m\ge 2$ this is $(m+1)x+2$.

Lean: \texttt{JacobiSum.jacobi} (the DLMF~18.5.7 sum), \texttt{tj} (its
$k$-th summand in $u=\tfrac{x-1}{2}$), the ratio identities
\texttt{tj\_succ}, \texttt{tj\_pred\_v}, \texttt{tj\_pred\_u} and their
combinations \texttt{tj\_pred}, \texttt{tj\_pred2}, the telescoping
certificate \texttt{W} with \texttt{summand\_telescopes}, \texttt{W\_zero},
\texttt{W\_top}, and DLMF~18.9.2 itself, \texttt{jacobi\_three\_term};
\texttt{Recurrence.jacobi\_zero}, \texttt{jacobi\_one}, \texttt{c0\_eq},
\texttt{den\_pos}, \texttt{recurrence\_form}, \texttt{march},
\texttt{march\_eq\_jacobi}.  No \texttt{sorry} in the chain.

\emph{Proof of DLMF~18.9.2 for the sum (as formalised).}  Let
$t_k=\binom{n+\alpha}{n-k}\binom{n+\beta}{k}u^k(u+1)^{n-k}$ for $0\le k\le n$
and $t_k=0$ otherwise.  From $\binom{N}{j+1}(j+1)=\binom{N}{j}(N-j)$ and
$\binom{N}{j}(N+1)=\binom{N+1}{j}(N+1-j)$:
\begin{align*}
(\alpha+k+1)(k+1)(u+1)\,t_{k+1}&=(n-k)(n+\beta-k)\,u\,t_k,\\
(n+1+\alpha)(n+1+\beta)(u+1)\,t^{n}_k&=(n+1-k)(n+1+\beta-k)\,t^{n+1}_k,\\
(n+1+\alpha)(n+1+\beta)\,u\,t^{n}_k&=(\alpha+k+1)(k+1)\,t^{n+1}_{k+1},
\end{align*}
and the difference of the last two expresses $t^n_k$ through $t^{n+1}_k$,
$t^{n+1}_{k+1}$ with no division.  Applying it twice writes the recurrence's
summand at degree $N$, times $D'=(N+\alpha)(N+\beta)(N+\alpha-1)(N+\beta-1)$,
through $t_k,t_{k+1},t_{k+2}$ alone, and it equals $W(k+1)-W(k)$ for
$W(k)=-q_0(k)\,t_k+\bigl(q_2(k-1)-C(u+1)(k+1)(k+1+\alpha)\bigr)t_{k+1}$
(explicit polynomials $q_0,q_2,C$, found with sympy and checked by
\texttt{ring}); $W(0)=0$ by the first identity at $k=0$ and $W(N+1)=0$
because $t$ vanishes past $N$, so the sum telescopes to zero.

\section{Exponent carry}
\label{sec:carry}

Every lane stores $v\cdot 2^{ex}$ with $v$ float32 and $ex\in\mathbb{Z}$.
Initially $v=2^{f-\lfloor f\rfloor}\in[1,2)$ and $ex=\lfloor f\rfloor$, where
\begin{equation}
\label{eq:f}
f=\alpha\log_2\sin\tfrac{\theta}{2}+\beta\log_2\cos\tfrac{\theta}{2}.
\end{equation}
Whenever $\max(\lvert v_\ell\rvert,\lvert v_{\ell-1}\rvert)$ leaves
$[2^{-24},2^{24}]$, both lanes are multiplied by $2^{\mp 24}$ and $ex$
moves by $\pm 24$.

\begin{lemma}
\label{lem:rescale}
The rescaling is exact and commutes with the recurrence.
\end{lemma}

\begin{proof}
Multiplying a binary float by a power of two changes only the exponent
field, hence is exact in the absence of overflow or underflow.  With
$\lvert v\rvert\le 2^{24}$ times the growth of one step, and the float32
range $2^{\pm 126}$, overflow before the guard would require growth
$>2^{100}$.  From~\eqref{eq:c} one has
$\lvert c_1\rvert,\lvert c_0\rvert,\lvert c_b\rvert=O(m)$, so growth per
step is $\le 2^{15}$ at $m<2^{14}$.  The recurrence is linear and
homogeneous in $(v_\ell,v_{\ell-1})$, so a common scale persists.
\end{proof}

The band $2^{-24}$ is chosen so that $v\cdot r$ with
$\lvert r\rvert$ of class $2^{24}$ and a $256$-lane sum stays inside
float32.

Lean: \texttt{ExponentCarry.lane\_rescale}, \texttt{guard\_large},
\texttt{guard\_small}, \texttt{step\_homogeneous}, \texttt{rec2\_scaled} /
\texttt{march\_eq\_rec2} / \texttt{march\_scaled} (the march from seeds
scaled by $\lambda$ is $\lambda$ times the march at every degree),
\texttt{no\_overflow\_before\_guard}.

\section{Emit}
\label{sec:emit}

At degree $\ell$ the program holds $(v_\theta,ex_\theta)$ on $256$ lanes and
reduces $\sum_\theta v_\theta\,2^{ex_\theta-\emax}\,r_c(\theta)$ with
$\emax=\max_\theta ex_\theta$.  Lanes with $ex_\theta-\emax<-126$ are flushed
to zero.

\begin{lemma}[flush bound]
\label{lem:flush}
Flushed lanes change the tile partial by less than $2^{-70}$ relative to
the largest lane's term.
\end{lemma}

\begin{proof}
Stored $\lvert v\rvert$ lie in $[2^{-24},2^{24}]$.  The largest term is at
least $2^{-24}\cdot 2^{\emax}\cdot\lvert r\rvert$; a flushed term is at most
$2^{24}\cdot 2^{\emax-127}\cdot\lvert r\rvert$.  The ratio is $2^{-79}$, and
there are at most $2^8$ lanes.
\end{proof}

From~\eqref{eq:d} the tile partial of~\eqref{eq:A} is
\begin{equation}
\label{eq:tile}
A_c^{\mathrm{tile}}(\ell,m)
  =(-1)^{m+s}\,
   2^{\varepsilon\,\LG+\emax}
   \sum_\theta
   \bigl[v_\theta\, 2^{ex_\theta-\emax}\bigr]\, r_c(\theta).
\end{equation}
Write $\varepsilon\,\LG=\mathrm{lex}+\mathrm{frac}$ with
$\mathrm{lex}=\lfloor\varepsilon\,\LG\rfloor$ and
$\mathrm{frac}\in[0,1)$.  The kernel stores
\begin{equation}
\label{eq:p}
p=\fl{32}\Bigl(\fl{32}\bigl(\textstyle\sum_\theta\mathrm{val}_\theta\, r_c(\theta)\bigr)
   \cdot\fl{32}\bigl((-1)^{m+s}\,2^{\mathrm{frac}}\bigr)\Bigr),
\qquad
e=\mathrm{lex}+\emax,
\end{equation}
and the driver returns $\fl{64}(p)\cdot 2^e$, summing tiles in float64.

\begin{proposition}
\label{prop:binade}
$\fl{64}(p)\,2^e$ equals~\eqref{eq:tile} with relative error at most
$(1+u)^2(1+\delta)-1$, where $u=2^{-24}$ and $\delta$ is the relative error
of the float32 tree sum.  Multiplication by $2^e$ is exact for
$-1022\le e\le 1023$.  A partial with $e<-1022$ is below $2^{-1022}$ in
magnitude and is dropped.
\end{proposition}

\begin{proof}
$2^e$ is exactly representable in float64 on that range, and multiplication
by it only changes the exponent field.  The two float32 roundings are
$\fl{32}(2^{\mathrm{frac}})$ and the product; each contributes a factor in
$[1-u,1+u]$.  The tree error $\delta$ is the same term the previous
float64-emit kernel carried.  Harmonic coefficients are $O(1)$ sums of
these partials, so a contribution below $2^{-1022}$ is below float64
resolution.
\end{proof}

Measured relative change against the float64-emit kernel, over every lane:
$7.5\times 10^{-8}$, $7.3\times 10^{-8}$, $6.6\times 10^{-8}$ at
$N_{\mathrm{side}}=1024$, $2048$, $4096$, i.e.\ $\sim 2u$ as predicted.

Lean: \texttt{Emit.flushed\_lane\_small}, \texttt{flush\_bound},
\texttt{binade\_split}, \texttt{frac\_factor\_bounds},
\texttt{int\_binade\_is\_zpow}, \texttt{two\_roundings},
\texttt{dropped\_partial\_small}.

\section{Hemisphere fold}
\label{sec:fold}

\begin{lemma}
\label{lem:fold}
$d^\ell_{m,0}(\pi-\theta)=(-1)^{\ell+m}\, d^\ell_{m,0}(\theta)$.
\end{lemma}

\begin{proof}
$d^\ell_{m,0}(\theta)=\sqrt{(\ell-m)!/(\ell+m)!}\,P_\ell^m(\cos\theta)$ and
$P_\ell^m(-x)=(-1)^{\ell+m}P_\ell^m(x)$ (DLMF~14.7.17).  Equivalently,
$P_n^{(\alpha,\alpha)}(-x)=(-1)^n P_n^{(\alpha,\alpha)}(x)$ from the
explicit sum, and $\alpha=\beta=m$ at spin~$0$.
\end{proof}

HEALPix rings satisfy $\theta_{N-1-i}=\pi-\theta_i$ and share weight and
$\phi_0$.  The sum~\eqref{eq:A} therefore collapses, per northern lane, to
\begin{equation}
\label{eq:fold}
\sum_{\mathrm{north}} d_i\bigl[G_i+(-1)^{\ell+m}G_i'\bigr].
\end{equation}
One recurrence, two right-hand sides; the parity sign is applied outside
the kernel.  This halves the lane count for spin~$0$.

For spin~$2$ the code uses
$T[m,\pi-\theta,\ell]=(-1)^{\ell-m'}T[-m,\theta,\ell]$ with $m'=-s$
(measured residual $2.3\times 10^{-12}$).  It is a consequence of the general
reflection $P_n^{(\alpha,\beta)}(-x)=(-1)^nP_n^{(\beta,\alpha)}(x)$, which swaps
the two half-angle powers in~\eqref{eq:d}: for $s\le m\le\ell$,
$d^\ell_{m,-s}(\pi-\theta)=(-1)^{\ell+m}\,d^\ell_{m,s}(\theta)$, and
$d^\ell_{-m,-s}=(-1)^{m+s}d^\ell_{m,s}$ is read off the same Jacobi form.
Lean: \texttt{Mirror.jacobi\_reflect}, \texttt{rowNeg\_reflect},
\texttt{mirror\_identity}.

Lean: \texttt{Fold.jacobi\_symmetric\_parity}, \texttt{row0\_reflect},
\texttt{sum\_pairs\_even}, \texttt{sum\_pairs\_odd}, \texttt{fold\_even}.

\section{The arithmetic wall}
\label{sec:wall}

\begin{lemma}
\label{lem:live}
$\sum_{m<L}(L-m)=L(L+1)/2$.
\end{lemma}

Let $N$ be the number of $(m,\ell,\theta)$ triples visited.  Before polar
skip, $N=\tfrac12 L(L+1)\lceil n_\theta/2\rceil$ on the folded spin-$0$
route and $N=\tfrac12 L(L+1)\,n_\theta$ at spin~$2$.

\begin{center}
\small
\begin{tabular}{@{}rrrrr@{}}
\toprule
$N_{\mathrm{side}}$ & $L$ & rings (spin $0$ / $2$) & $N$ spin~$0$ & $N$ spin~$2$ \\
\midrule
$1024$ & $3072$  & $2048$ / $4095$  & $9.7\times 10^{9}$  & $1.9\times 10^{10}$ \\
$2048$ & $6144$  & $4096$ / $8191$  & $7.7\times 10^{10}$ & $1.5\times 10^{11}$ \\
$4096$ & $12288$ & $8192$ / $16383$ & $6.2\times 10^{11}$ & $1.2\times 10^{12}$ \\
\bottomrule
\end{tabular}
\end{center}

An analysis pass with $n_c$ channels costs at least $(2+n_c)N$ fused
multiply-adds.  On this card the measured float64 FMA rate is
$R_{64}=0.45\times 10^{12}\,\mathrm{s}^{-1}$ (datasheet
$0.94\times 10^{12}$).  Hence any float64 SIMT spin-$0$ analysis at
$N_{\mathrm{side}}=4096$ needs
\begin{equation}
\label{eq:T}
T\ge\frac{6\cdot 6.2\times 10^{11}}{0.94\times 10^{12}}\,\mathrm{s}
   =4.0\,\mathrm{s}
\end{equation}
at the datasheet rate ($8.3\,\mathrm{s}$ measured), against ducc0's
$2.02\,\mathrm{s}$ on the host.  No float64 kernel on this card can beat
the CPU reference at the top of the board.

Lean: \texttt{Wall.live\_degrees}, \texttt{time\_lower\_bound}.

\section{Polar skip}
\label{sec:polar}

For $m>(\ell+\tfrac12)\sin\theta$ the function $d^\ell_{m,-s}(\theta)$ is
in its forbidden region.  Writing $u=\sqrt{\sin\theta}\,d^\ell_{m,0}$,
\begin{equation}
\label{eq:legendre}
u''+\Bigl[(\ell+\tfrac12)^2-\frac{m^2-\tfrac14}{\sin^2\theta}\Bigr]u=0.
\end{equation}
With $\ell'=\ell+\tfrac12$ the WKB action from the turning point
$\theta_t$ ($\sin\theta_t=m/\ell'$) down to $\theta<\theta_t$ is
\begin{equation}
\label{eq:S}
S=\int_\theta^{\theta_t}
   \sqrt{\frac{m^2}{\sin^2 t}-{\ell'}^2}\,\mathrm{d}t
  \;\sim\;
  \frac{2\sqrt{2}}{3}\,
  \frac{\delta^{3/2}}{\sqrt{\ell'\sin\theta}\,\cos\theta},
\end{equation}
where $\delta=m-\ell'\sin\theta$.  This estimate is not formalised.

ducc0 neglects, on the ring at $\theta$, every order above
\begin{equation}
\label{eq:mlim}
\mlim
  =s\lvert\cos\theta\rvert
   +\sqrt{(\ell_{\max}\sin\theta+\ofs)^2-s^2\sin^2\theta},
\qquad
\ofs=\max(100,\,0.01\,\ell_{\max}).
\end{equation}

\begin{lemma}
\label{lem:mlim}
\eqref{eq:mlim} is the larger root of
$m^2-2s\lvert\cos\theta\rvert\, m+s^2-t_1^2=0$ on a ring with
$\cos^2\theta+\sin^2\theta=1$ and $t_1=\ell_{\max}\sin\theta+\ofs$.
For $s=0$ one has $\mlim=\ell_{\max}\sin\theta+\ofs$.
\end{lemma}

Every $\ell\le\ell_{\max}$ then has $\delta\ge\ofs$.  At
$\ell_{\max}=12287$, $\sin\theta=\tfrac12$, one finds $S\sim 19$, so
neglected values sit below $\sim 10^{-8}$ times the oscillatory amplitude.
The march skips a tile when $m>\max_{\mathrm{tile}}\mlim(\theta)$.  Because
the neglected set is the reference's own, agreement with ducc0 is unchanged
to the printed digit.

\begin{center}
\small
\begin{tabular}{@{}rrrr@{}}
\toprule
$N_{\mathrm{side}}$ & analysis (tile $256$) & synth.\ spin~$0$ ($1024$) & synth.\ spin~$2$ ($512$) \\
\midrule
$1024$ & $13.9\%$ & $2.5\%$  & $9.3\%$  \\
$2048$ & $17.5\%$ & $10.0\%$ & $14.8\%$ \\
$4096$ & $19.3\%$ & $15.1\%$ & $17.9\%$ \\
\bottomrule
\end{tabular}
\end{center}

Lean: \texttt{PolarSkip.mlim\_is\_root}, \texttt{mlim\_ge},
\texttt{mlim\_spin\_zero}, \texttt{skipped\_degrees}.

\appendix
\section{Lean correspondence}
\label{app:lean}

Build: \texttt{cd formal/lean \&\& $\sim$/.elan/bin/lake build}.
There is no \texttt{sorry}; \texttt{lake env lean CheckAxioms.lean} prints the
axioms of every main theorem (\texttt{propext}, \texttt{Classical.choice},
\texttt{Quot.sound}).

\begin{center}
\small
\begin{tabular}{@{}llp{7.4cm}@{}}
\toprule
Section & File & Content \\
\midrule
\S\ref{sec:closed}
  & \texttt{Indices.lean}
  & $\alpha+\beta=2M$, $\alpha^2-\beta^2=4ms$, both factorial branches, sign \\
\S\ref{sec:closed}
  & \texttt{WignerJacobi.lean}
  & Wigner's explicit sum equals the Jacobi closed form, both branches \\
\S\ref{sec:rec}
  & \texttt{JacobiSum.lean}
  & Jacobi sum; DLMF~18.9.2 by summand ratios and a telescoping certificate \\
\S\ref{sec:rec}
  & \texttt{Recurrence.lean}
  & seeds, rearrangement, $\mathrm{march}=\mathrm{jacobi}$ \\
\S\ref{sec:carry}
  & \texttt{ExponentCarry.lean}
  & exact rescaling, homogeneous step, no overflow before the guard \\
\S\ref{sec:emit}
  & \texttt{Emit.lean}
  & flush $2^{-71}$, binade split, $[1,2)$ mantissa, two roundings \\
\S\ref{sec:fold}
  & \texttt{Fold.lean}
  & $P_n^{(\alpha,\alpha)}(-x)=(-1)^n P_n^{(\alpha,\alpha)}(x)$, paired rings \\
\S\ref{sec:fold}
  & \texttt{Mirror.lean}
  & $P_n^{(\alpha,\beta)}(-x)=(-1)^nP_n^{(\beta,\alpha)}(x)$; the spin-$2$ mirror identity \\
\S\ref{sec:wall}
  & \texttt{Wall.lean}
  & Lemma~\ref{lem:live}; time bound from an FMA count and a rate \\
\S\ref{sec:polar}
  & \texttt{PolarSkip.lean}
  & Lemma~\ref{lem:mlim}; degrees removed by a cutoff \\
\bottomrule
\end{tabular}
\end{center}

Not formalised: the WKB action \eqref{eq:S}.

\begin{thebibliography}{9}
\bibitem{Edm57}
A.~R. Edmonds, \emph{Angular Momentum in Quantum Mechanics},
Princeton University Press, 1957.
\bibitem{DLMF}
NIST Digital Library of Mathematical Functions,
\url{https://dlmf.nist.gov/}.
\end{thebibliography}

\end{document}
